
dont ils sont déduits est compatible avec les structures d'anneaux, et grâce à (7.5.1), on trouve que l'on a
@@M0@@
en utilisant les accouplements @@M1@@ de 4.1 entre @@M2@@ et @@M3@@. De ceci on déduit formellement que \emph{si} @@M4@@ \emph{est un élément de} @@M5@@ (\emph{resp}. \emph{de} @@M6@@) \emph{tel que} @@M7@@ (\emph{resp}. @@M8@@) \emph{soit numériquement équivalent à zéro} (4.5.1), \emph{alors il en est de même de} @@M9@@. La réciproque est vraie si (7.9.2) (ou ce qui revient au même, l'un ou l'autre homomorphisme (7.9.1) qui s'en déduit par passage aux gradués associés) est un épimorphisme mod.torsion, - circonstance qui n'est cependant réalisée qu'assez exceptionnellement.

\textbf{7.9.4.} Il est plausible, du moins lorsque @@M10@@ est propre et \emph{lisse}, que la réciproque envisagée soit toujours valable, i.e. que @@M11@@ (resp. @@M12@@) soit numériquement équivalent à zéro, dès qu'il en est ainsi de @@M13@@, ce qui impliquerait que les homomorphismes (7.9.1) induisent des \emph{monomorphismes} pour les groupes quotients correspondants par la relation d'équivalence numérique @@M14@@ et @@M15@@. Du moins ce qui précède implique que le groupe @@M16@@ est isomorphe à un quotient d'un sous-groupe de @@M17@@ (et en particulier, il a un rang au plus égal à celui de @@M18@@), et de même pour @@M19@@, et aussi pour @@M20@@ lui-même.

\subsection*{7.10. Compatibilité avec les images directes et images inverses}
Donnons-nous un morphisme
@@M21@@
de schémas noethériens et plats sur @@M22@@. On désigne par @@M23@@ (resp. @@M24@@) l'homomorphisme de spécialisation pour @@M25@@ (resp. @@M26@@), et par @@M27@@ (resp. @@M28@@)
